\section*{Capítulo \ref{ch:b2:liquid-vapour-diagrams} --- Diagramas binários líquido--vapor e destilação}

\begin{solution}{exo:b2:liquid-vapour-diagrams:1}
Benzeno \qty{80.1}{\degreeCelsius}, tolueno \qty{110.6}{\degreeCelsius}. O líquido
de composição 0.2 começa a ferver a cerca de \qty{102}{\degreeCelsius}; seu primeiro
vapor tem composição 0.38.
\end{solution}

\begin{solution}{exo:b2:liquid-vapour-diagrams:2}
$p = 0.5 \times 1.010 + 0.5 \times 0.388 = \qty{0.699}{bar}$;
$y = 0.505/0.699 = 0.722$.
\end{solution}

\begin{solution}{exo:b2:liquid-vapour-diagrams:3}
$n_V/n = (0.30 - 0.20)/(0.45 - 0.20) = 0.40$: \qty{4.0}{mol} de vapor e
\qty{6.0}{mol} de líquido. Verificação: $6.0 \times 0.20 + 4.0 \times 0.45 = 3.0 = 10
\times 0.30$.
\end{solution}

\begin{solution}{exo:b2:liquid-vapour-diagrams:4}
(a) $v = 2 + 2 - 1 = 3$: $T$, $p$ e $x$ são livres. (b) $v = 2 + 2 - 2 = 2$; a
pressão fixa, 1. (c) O \omterm{def:b2:liquid-vapour-diagrams:azeotrope}{azeótropo} acrescenta a relação $x = y$: $v = 1$; a pressão
fixa, 0 --- o \omterm{def:b2:liquid-vapour-diagrams:azeotrope}{azeótropo} ferve a uma temperatura fixa, como um líquido puro,
embora sua composição mude com a pressão.
\end{solution}

\begin{solution}{exo:b2:liquid-vapour-diagrams:5}
No \omterm{def:b2:liquid-vapour-diagrams:bubble-dew}{ponto de orvalho}, $yp/p_1^* + (1 - y)p/p_2^* = 1$. A \qty{100}{\degreeCelsius}:
$0.3 \times 1.013/1.80 + 0.7 \times 1.013/0.742 = 1.124 > 1$ (frio demais); a
\qty{105}{\degreeCelsius}: $0.148 + 0.824 = 0.971 < 1$. Interpolando, $T \approx
100 + 5 \times 0.124/0.153 = \qty{104}{\degreeCelsius}$. A primeira gota tem $x =
yp/p_1^* \approx 0.3 \times 1.013/2.0 = 0.15$.
\end{solution}

\begin{solution}{exo:b2:liquid-vapour-diagrams:6}
As primeiras gotas têm a composição do vapor sobre o líquido em ebulição,
$y = 0.71$, a \qty{92.1}{\degreeCelsius}. Como o benzeno sai preferencialmente, o
líquido fica mais pobre nele, seu ponto de ebulição sobe e o destilado também
fica mais pobre; as últimas gotas são quase tolueno puro. A destilação simples
é um único estágio de equilíbrio, repetido sobre um líquido que não para de
mudar: cada fração é uma mistura.
\end{solution}

\begin{solution}{exo:b2:liquid-vapour-diagrams:7}
Alimentação 0.05, do lado da água em relação ao \omterm{def:b2:liquid-vapour-diagrams:azeotrope}{azeótropo}: o topo fornece o
\omterm{def:b2:liquid-vapour-diagrams:azeotrope}{azeótropo} (0.88 em mols, 95~\% em massa), o refervedor retém a água.
Alimentação 0.95, do lado do etanol: o topo ainda fornece o \omterm{def:b2:liquid-vapour-diagrams:azeotrope}{azeótropo}, o ponto
de ebulição mais baixo, e o refervedor retém o etanol puro.
\end{solution}

\begin{solution}{exo:b2:liquid-vapour-diagrams:8}
A mistura ferve quando $0.10 + p^*_{\text{água}} = 1.013$, isto é,
$p^*_{\text{água}} = \qty{0.91}{bar}$, perto de \qty{97}{\degreeCelsius}. Razão de
massas $0.10 \times 136/(0.91 \times 18.02) = 0.83$: \qty{0.83}{g} de essência por
grama de água.
\end{solution}

\begin{solution}{exo:b2:liquid-vapour-diagrams:9}
À temperatura ambiente, duas camadas L$_2$ e L$_1$ (a composição 0.4 está na
lacuna). No aquecimento, as camadas fervem juntas na temperatura do
\omterm{def:b2:liquid-vapour-diagrams:miscibility-gap}{heteroazeótropo}, que permanece constante; o primeiro vapor tem a composição
heteroazeotrópica (o ponto). Como 0.4 fica do lado L$_2$ do ponto, a camada
L$_1$ se esgota primeiro; a temperatura então sobe, a camada L$_2$ restante
fica mais pobre em 1 ao longo de sua \omterm{def:b2:liquid-vapour-diagrams:bubble-dew}{curva de bolha} e o vapor segue a \omterm{def:b2:liquid-vapour-diagrams:bubble-dew}{curva de
orvalho}, até que o ponto atinja a \omterm{def:b2:liquid-vapour-diagrams:bubble-dew}{curva de orvalho}: a última gota evapora.
\end{solution}

\begin{solution}{exo:b2:liquid-vapour-diagrams:10}
$N_{\min} = \ln(99 \times 99)/\ln 2.47 = 9.19/0.904 = 10.2$, contra 6.5: passar
de 95~\% a 99~\% de pureza nas duas extremidades custa mais da metade dos
pratos a mais. Cada fator adicional em $x/(1 - x)$ custa o mesmo número de
pratos, e a pureza é medida pela impureza restante, que só cai geometricamente.
\end{solution}

\begin{solution}{exo:b2:liquid-vapour-diagrams:11}
O \omterm{def:b2:liquid-vapour-diagrams:binary-diagram}{diagrama isobárico} sozinho não o fornece: são necessárias $p_1^*$ e $p_2^*$ a
\qty{80}{\degreeCelsius}, dadas pelas equações de Antoine (1.010 e
\qty{0.388}{bar}). Então $p = p_2^* + x(p_1^* - p_2^*)$ e $p = 1/(y/p_1^* + (1 -
y)/p_2^*)$. O líquido é estável a alta pressão (comprimir um vapor o
condensa), e a baixa temperatura no \omterm{def:b2:liquid-vapour-diagrams:binary-diagram}{diagrama isobárico}: os dois diagramas
são invertidos um em relação ao outro.
\end{solution}

\begin{solution}{exo:b2:liquid-vapour-diagrams:12}
Como $y$ cresce com $x$ e $y(1 - y) > 0$, $dp/dx$ tem o sinal de $y - x$:
a pressão cresce com $x$ onde o vapor é mais rico em 1. Para uma \omterm{def:b2:chemical-potential:ideal-mixture}{mistura
ideal}, $y > x$ para todo $0 < x < 1$ (\cref{prop:b2:liquid-vapour-diagrams:richer}):
$p$ é estritamente monótona, não tem extremo, e não há \omterm{def:b2:liquid-vapour-diagrams:azeotrope}{azeótropo}.
\end{solution}

\begin{solution}{pb:b2:liquid-vapour-diagrams:1}
\textbf{1.} $T = B/(A - \log_{10}1.013) - C$: benzeno $1203.835/4.01253 + 53.226 =
\qty{353.2}{K}$, \qty{80.1}{\degreeCelsius}; tolueno $1343.943/4.07266 + 53.773 =
\qty{383.8}{K}$, \qty{110.6}{\degreeCelsius}.
\textbf{2.} A \qty{363.15}{K}: $p_1^* = \qty{1.361}{bar}$, $p_2^* = \qty{0.542}{bar}$.
\textbf{3.} $x = (1.013 - 0.542)/(1.361 - 0.542) = 0.575$; $y = 0.575 \times
1.361/1.013 = 0.773$.
\textbf{4.} $x = (1.013 - 0.742)/(1.800 - 0.742) = 0.256$; $y = 0.256 \times
1.800/1.013 = 0.455$.
\textbf{5.} \omterm{def:b2:liquid-vapour-diagrams:bubble-dew}{Pontos de bolha} (0, 110.6), (0.256, 100), (0.575, 90), (1, 80.1);
\omterm{def:b2:liquid-vapour-diagrams:bubble-dew}{pontos de orvalho} (0, 110.6), (0.455, 100), (0.773, 90), (1, 80.1) --- o fuso
da figura.
\textbf{6.} $\frac12(p_1^* + p_2^*) = 1.013$: a \qty{90}{\degreeCelsius} 0.952, a
\qty{95}{\degreeCelsius} 1.102; interpolando, $90 + 5 \times 0.061/0.150 \approx
\qty{92}{\degreeCelsius}$ (\qty{92.1}{\degreeCelsius} exatamente).
\textbf{7.} $v = 2$, 1 a pressão fixa: cada temperatura fixa $x$ e $y$,
e é por isso que o diagrama é formado por duas curvas.
\textbf{8.} $x = (1.013 - 0.636)/(1.569 - 0.636) = 0.404$; $y = 0.404 \times
1.569/1.013 = 0.626$.
\textbf{9.} $n_V = 100 \times (0.500 - 0.404)/(0.626 - 0.404) = \qty{43}{mol}$,
$n_L = \qty{57}{mol}$.
\textbf{10.} Vapor $43 \times 0.626 = \qty{26.9}{mol}$; líquido $57 \times 0.404
= \qty{23.0}{mol}$; total \qty{49.9}{mol} $\approx 50$.
\textbf{11.} $26.9/50 = 54~\%$ do benzeno, em um vapor de apenas 63~\%.
\textbf{12.} A \qty{100}{\degreeCelsius}, o vapor em equilíbrio tem $y =
0.455 < 0.5$: o ponto da alimentação fica acima da \omterm{def:b2:liquid-vapour-diagrams:bubble-dew}{curva de orvalho}, a
alimentação está inteiramente vaporizada e nada é separado.
\textbf{13.} Em seu \omterm{def:b2:liquid-vapour-diagrams:bubble-dew}{ponto de orvalho}, \qty{98.8}{\degreeCelsius} (pelo método do
\cref{exo:b2:liquid-vapour-diagrams:5}).
\textbf{14.} $\alpha = p_1^*/p_2^*$; a \qty{80.1}{\degreeCelsius}, $1.013/0.390 =
2.60$.
\textbf{15.} A \qty{110.6}{\degreeCelsius}, $2.377/1.013 = 2.35$.
\textbf{16.} $\sqrt{2.60 \times 2.35} = 2.47$.
\textbf{17.} $y = xp_1^*/p$ e $1 - y = (1 - x)p_2^*/p$; divida.
\textbf{18.} Nada é retirado: entre dois pratos, o vapor que sobe e o líquido
que desce transportam a mesma quantidade de matéria e de cada componente,
$V y_k = L x_{k+1}$ com $V = L$; logo $x_{k+1} = y_k$.
\textbf{19.} Com $r = x/(1 - x)$: $r(y_k) = \alpha\,r(x_k)$ e $x_{k+1} = y_k$; logo
$r(x_{k+1}) = \alpha\,r(x_k)$ e $r(x_D) = \alpha^N r(x_B)$; $N = \ln[r(x_D)/
r(x_B)]/\ln\alpha$.
\textbf{20.} $N_{\min} = \ln(19 \times 19)/\ln 2.47 = 5.889/0.904 = 6.5$.
\textbf{21.} Uma coluna precisa fornecer produto; por isso opera a uma \omterm{def:b2:liquid-vapour-diagrams:fractional-distillation}{razão de
refluxo} finita, que exige mais pratos; e um prato real não atinge o equilíbrio
(sua eficiência é menor que um).
\textbf{22.} Sete degraus, o último terminando em 0.034, abaixo de 0.05: 6.5
pratos em contagem contínua, sete inteiros (contando o refervedor como um).
\textbf{23.} $x = (95/46.07)/(95/46.07 + 5/18.02) = 0.881$.
\textbf{24.} No topo, no melhor dos casos, o \omterm{def:b2:liquid-vapour-diagrams:azeotrope}{azeótropo} (0.88 em mols); no fundo,
a água. Cinquenta pratos não atravessam o \omterm{def:b2:liquid-vapour-diagrams:azeotrope}{azeótropo}.
\textbf{25.} A refluxo total, $\boldsymbol{N_{\min} \approx 6.5}$ \omterm{def:b2:liquid-vapour-diagrams:fractional-distillation}{pratos
teóricos}.
plates.
\end{solution}
